\section{Conclusion}
\label{sec:conclusion}

This paper examines a post-2022 age-exposure gradient in Norwegian employment using aggregated cell counts from population administrative data (2021--2026). In the private sector, the most-exposed quintile declines relative to the least-exposed among workers aged 41 to 50 rather than among the youngest workers, and there is no corresponding wage divergence. Public-sector employment shows no gradient.

AI-related adjustment, the post-COVID correction in technology hiring, monetary tightening by Norges Bank, and seniority-based retention could each generate the same age-exposure gradient, and the aggregate data we have cannot tell them apart. \citet{humlum2026} document the same aggregate pattern in Denmark and show that observable workplace adoption of AI chatbots accounts for at most a small fraction of it. Worker- and firm-level designs in \citet{brynjolfsson2025}, \citet{lodefalk2026}, and \citet{humlum2026} answer related questions at finer units of observation.
